% ================================================================
%  SESIÓN 06 — Figuras planas en la ciudad
%  Almería en Cifras y Formas · Matemáticas 1.º ESO
%  IES Al-Ándalus · 2025-2026
% ================================================================
\documentclass[aspectratio=169, 12pt, spanish]{beamer}
\usepackage{almeria-beamer}

\renewcommand{\SessionNumber}{06}
\renewcommand{\SessionTitle}{Figuras planas en la ciudad}
\renewcommand{\SessionName}{Sesión 06}
\renewcommand{\SessionObjective}{%
  Identificar y clasificar figuras planas (polígonos) en elementos
  arquitectónicos de Almería, reconociendo sus propiedades básicas
  (lados, vértices, ángulos).%
}
\renewcommand{\BlockName}{Bloque 2: Geometría urbana}

\begin{document}

% ---------------------------------------------------------------
%  PORTADA
% ---------------------------------------------------------------
\begin{frame}[plain]
  \titlepage
\end{frame}

% ---------------------------------------------------------------
%  ÍNDICE
% ---------------------------------------------------------------
\begin{frame}{Contenidos de hoy}
  \begin{columns}[t]
    \begin{column}{0.48\textwidth}
      \begin{itemize}
        \item Figura plana y polígono
        \item Clasificación de polígonos
        \item Propiedades: lados, vértices, ángulos
        \item Polígonos regulares e irregulares
        \item Cuadriláteros especiales
      \end{itemize}
    \end{column}
    \begin{column}{0.48\textwidth}
      \begin{itemize}
        \item Tipos de triángulos
        \item Suma de ángulos interiores
        \item Polígonos en Almería
        \item Ejercicios Bloom
        \item Evidencia del día
      \end{itemize}
    \end{column}
  \end{columns}
  \vspace{0.4cm}
  \begin{notabox}[Niveles de Bloom]
    \bloomtag{Recordar}\quad
    \bloomtag{Comprender}\quad
    \bloomtag{Aplicar}\quad
    \bloomtag{Analizar}
  \end{notabox}
\end{frame}

% ---------------------------------------------------------------
\seccionframe{Figuras planas y polígonos}
% ---------------------------------------------------------------

% ---------------------------------------------------------------
%  FRAME 3 — Figura plana y polígono
% ---------------------------------------------------------------
\begin{frame}{¿Qué es una figura plana?}
  \begin{columns}[c]
    \begin{column}{0.54\textwidth}
      \begin{definicionbox}{Figura plana}
        Figura geométrica que se puede dibujar completamente en
        un plano (en dos dimensiones: largo y ancho).
        No tiene profundidad.
      \end{definicionbox}
      \vspace{0.3cm}
      \begin{definicionbox}{Polígono}
        Figura plana \textbf{cerrada} formada por segmentos de recta
        (llamados \textbf{lados}). Los segmentos se unen en puntos
        llamados \textbf{vértices}.
      \end{definicionbox}
    \end{column}
    \begin{column}{0.42\textwidth}
      \begin{center}
      \begin{tikzpicture}[scale=0.85]
        % polígono ejemplo (pentágono irregular)
        \draw[zona] (0,0) -- (2,0) -- (2.6,1.4) -- (1.3,2.2) -- (-0.4,1.3) -- cycle;
        \foreach \p/\lbl/\pos in {
          (0,0)/A/below left,
          (2,0)/B/below right,
          (2.6,1.4)/C/right,
          (1.3,2.2)/D/above,
          (-0.4,1.3)/E/left}{
          \node[punto] at \p {};
          \node[\pos, font=\bfseries\scriptsize, color=AlmAzulM] at \p {\lbl};
        }
        % etiqueta lados
        \node[color=AlmTerra, font=\tiny\bfseries] at (1,-0.3) {lado};
        % etiqueta vértice
        \draw[->, AlmTerra, thin] (0.6,-0.5) -- (-0.05,-0.05);
        \node[color=AlmTerra, font=\tiny\bfseries] at (0.9,-0.55) {vértice};
      \end{tikzpicture}
      \end{center}
      \begin{notabox}[Clave]
        \scriptsize Un círculo \textbf{NO} es un polígono:
        no tiene lados rectos.
      \end{notabox}
    \end{column}
  \end{columns}
\end{frame}

% ---------------------------------------------------------------
%  FRAME 4 — Galería de polígonos (TikZ)
% ---------------------------------------------------------------
\begin{frame}{Galería de polígonos: de 3 a 8 lados}
  \begin{center}
  \begin{tikzpicture}[scale=0.78]
    % ---- Triángulo ----
    \draw[zona] (0,0) -- (1.2,0) -- (0.6,1.0) -- cycle;
    \node[below, font=\tiny\bfseries, color=AlmAzulM] at (0.6,-0.12) {Triángulo};
    \node[below, font=\tiny, color=AlmGris] at (0.6,-0.38) {3 lados / 180°};
    % ---- Cuadrilátero ----
    \draw[zona] (1.7,0) -- (3.1,0) -- (3.1,1.0) -- (1.7,1.0) -- cycle;
    \node[below, font=\tiny\bfseries, color=AlmAzulM] at (2.4,-0.12) {Cuadrilátero};
    \node[below, font=\tiny, color=AlmGris] at (2.4,-0.38) {4 lados / 360°};
    % ---- Pentágono ----
    \begin{scope}[shift={(4.0,0.5)}]
      \draw[zona] \foreach \i in {0,1,2,3,4}{
        -- ({cos(90+\i*72)*0.6},{sin(90+\i*72)*0.6})
      } -- cycle;
    \end{scope}
    \node[below, font=\tiny\bfseries, color=AlmAzulM] at (4.0,-0.12) {Pentágono};
    \node[below, font=\tiny, color=AlmGris] at (4.0,-0.38) {5 lados / 540°};
    % ---- Hexágono ----
    \begin{scope}[shift={(5.7,0.5)}]
      \draw[zona] \foreach \i in {0,1,2,3,4,5}{
        -- ({cos(90+\i*60)*0.6},{sin(90+\i*60)*0.6})
      } -- cycle;
    \end{scope}
    \node[below, font=\tiny\bfseries, color=AlmAzulM] at (5.7,-0.12) {Hexágono};
    \node[below, font=\tiny, color=AlmGris] at (5.7,-0.38) {6 lados / 720°};
    % ---- Heptágono ----
    \begin{scope}[shift={(7.4,0.5)}]
      \draw[zona] \foreach \i in {0,1,2,3,4,5,6}{
        -- ({cos(90+\i*51.43)*0.6},{sin(90+\i*51.43)*0.6})
      } -- cycle;
    \end{scope}
    \node[below, font=\tiny\bfseries, color=AlmAzulM] at (7.4,-0.12) {Heptágono};
    \node[below, font=\tiny, color=AlmGris] at (7.4,-0.38) {7 lados / 900°};
    % ---- Octógono ----
    \begin{scope}[shift={(9.1,0.5)}]
      \draw[zona] \foreach \i in {0,1,2,3,4,5,6,7}{
        -- ({cos(90+\i*45)*0.6},{sin(90+\i*45)*0.6})
      } -- cycle;
    \end{scope}
    \node[below, font=\tiny\bfseries, color=AlmAzulM] at (9.1,-0.12) {Octógono};
    \node[below, font=\tiny, color=AlmGris] at (9.1,-0.38) {8 lados / 1080°};
  \end{tikzpicture}
  \end{center}
  \vspace{0.15cm}
  \begin{teoriabox}
    \small\textbf{Regla clave:} en todo polígono, el número de lados
    $=$ número de vértices $=$ número de ángulos interiores.
    La suma de ángulos interiores de un polígono de $n$ lados es
    $(n-2)\times 180°$.
  \end{teoriabox}
\end{frame}

% ---------------------------------------------------------------
%  FRAME 5 — Regular vs. Irregular
% ---------------------------------------------------------------
\begin{frame}{Polígono regular vs. irregular}
  \begin{columns}[c]
    \begin{column}{0.48\textwidth}
      \begin{definicionbox}{Polígono regular}
        Todos los lados son \textbf{iguales}
        \textbf{y} todos los ángulos son \textbf{iguales}.
        \\[4pt]
        Ejemplos: triángulo equilátero, cuadrado, hexágono regular.
      \end{definicionbox}
      \vspace{0.3cm}
      \begin{center}
      \begin{tikzpicture}[scale=0.8]
        \begin{scope}[shift={(0,0)}]
          \draw[rectov, fill=AlmVerde!15] \foreach \i in {0,1,2}{
            -- ({cos(90+\i*120)*0.7},{sin(90+\i*120)*0.7})
          } -- cycle;
          \node[below, font=\scriptsize\bfseries, color=AlmVerde] at (0,-0.8) {Regular};
        \end{scope}
        \begin{scope}[shift={(1.9,0)}]
          \draw[rectov, fill=AlmVerde!15] (0,0) -- (1.2,0) -- (1.2,1.2) -- (0,1.2) -- cycle;
          \node[below, font=\scriptsize\bfseries, color=AlmVerde] at (0.6,-0.12) {Regular};
        \end{scope}
      \end{tikzpicture}
      \end{center}
    \end{column}
    \begin{column}{0.48\textwidth}
      \begin{definicionbox}{Polígono irregular}
        Los lados \textbf{o} los ángulos son diferentes entre sí.
        \\[4pt]
        Muchas figuras de la arquitectura real son irregulares.
      \end{definicionbox}
      \vspace{0.3cm}
      \begin{center}
      \begin{tikzpicture}[scale=0.8]
        \begin{scope}[shift={(0,0)}]
          \draw[rectot, fill=AlmTerra!12]
            (0,0) -- (1.4,0) -- (1.0,0.8) -- (0.3,1.1) -- cycle;
          \node[below, font=\scriptsize\bfseries, color=AlmTerra] at (0.7,-0.12) {Irregular};
        \end{scope}
        \begin{scope}[shift={(2.2,0)}]
          \draw[rectot, fill=AlmTerra!12]
            (0,0.3) -- (1.1,0) -- (1.3,0.9) -- (0.8,1.3) -- (0.1,1.0) -- cycle;
          \node[below, font=\scriptsize\bfseries, color=AlmTerra] at (0.65,-0.12) {Irregular};
        \end{scope}
      \end{tikzpicture}
      \end{center}
    \end{column}
  \end{columns}
\end{frame}

% ---------------------------------------------------------------
%  FRAME 6 — Cuadriláteros especiales
% ---------------------------------------------------------------
\begin{frame}{Cuadriláteros especiales}
  \begin{center}
  \begin{tikzpicture}[scale=0.78, every node/.style={font=\tiny}]
    % Rectángulo
    \draw[zona] (0,0) rectangle (1.8,1.0);
    \node[below=2pt] at (0.9,-0.0) {\textbf{Rectángulo}};
    \node[below=14pt, color=AlmGris] at (0.9,0) {4 ángulos rectos};
    \node[below=24pt, color=AlmGris] at (0.9,0) {lados opuestos =};
    % Cuadrado
    \draw[zona] (2.4,0) rectangle (3.6,1.2);
    \node[below=2pt] at (3.0,-0.0) {\textbf{Cuadrado}};
    \node[below=14pt, color=AlmGris] at (3.0,0) {4 ángulos rectos};
    \node[below=24pt, color=AlmGris] at (3.0,0) {4 lados iguales};
    % Rombo
    \draw[zona] (4.6,0.5) -- (5.4,0) -- (6.2,0.5) -- (5.4,1.0) -- cycle;
    \node[below=2pt] at (5.4,-0.0) {\textbf{Rombo}};
    \node[below=14pt, color=AlmGris] at (5.4,0) {4 lados iguales};
    \node[below=24pt, color=AlmGris] at (5.4,0) {ángulos ≠ 90°};
    % Trapecio
    \draw[zona] (7.0,0) -- (8.6,0) -- (8.2,1.0) -- (7.4,1.0) -- cycle;
    \node[below=2pt] at (7.8,-0.0) {\textbf{Trapecio}};
    \node[below=14pt, color=AlmGris] at (7.8,0) {1 par de lados};
    \node[below=24pt, color=AlmGris] at (7.8,0) {paralelos};
    % Paralelogramo
    \draw[zona] (9.4,0) -- (11.0,0) -- (10.6,1.0) -- (9.0,1.0) -- cycle;
    \node[below=2pt] at (10.0,-0.0) {\textbf{Paralelogramo}};
    \node[below=14pt, color=AlmGris] at (10.0,0) {2 pares paralelos};
    \node[below=24pt, color=AlmGris] at (10.0,0) {lados opuestos =};
  \end{tikzpicture}
  \end{center}
  \vspace{0.1cm}
  \begin{notabox}[Jerarquía]
    \small Cuadrado $\subset$ Rectángulo $\subset$ Paralelogramo $\subset$ Cuadrilátero
  \end{notabox}
\end{frame}

% ---------------------------------------------------------------
%  FRAME 7 — Tipos de triángulos
% ---------------------------------------------------------------
\begin{frame}{Tipos de triángulos}
  \begin{columns}[t]
    \begin{column}{0.50\textwidth}
      \textbf{\color{AlmAzulM}Por sus lados:}
      \vspace{0.2cm}
      \begin{center}
      \begin{tikzpicture}[scale=0.75]
        % Equilátero
        \draw[rectov, fill=AlmVerde!12]
          (0,0) -- (1.0,0) -- (0.5,0.87) -- cycle;
        \node[below, font=\tiny\bfseries, color=AlmVerde] at (0.5,-0.05) {Equilátero};
        \node[below, font=\tiny, color=AlmGris] at (0.5,-0.25) {3 lados iguales};
        % Isósceles
        \draw[zona]
          (1.8,0) -- (3.2,0) -- (2.5,1.1) -- cycle;
        \node[below, font=\tiny\bfseries, color=AlmAzulM] at (2.5,-0.05) {Isósceles};
        \node[below, font=\tiny, color=AlmGris] at (2.5,-0.25) {2 lados iguales};
        % Escaleno
        \draw[rectot, fill=AlmTerra!10]
          (4.0,0) -- (5.6,0) -- (5.2,1.0) -- cycle;
        \node[below, font=\tiny\bfseries, color=AlmTerra] at (4.9,-0.05) {Escaleno};
        \node[below, font=\tiny, color=AlmGris] at (4.9,-0.25) {todos distintos};
      \end{tikzpicture}
      \end{center}
    \end{column}
    \begin{column}{0.46\textwidth}
      \textbf{\color{AlmAzulM}Por sus ángulos:}
      \vspace{0.2cm}
      \begin{center}
      \begin{tikzpicture}[scale=0.75]
        % Acutángulo
        \draw[rectov, fill=AlmVerde!12]
          (0,0) -- (1.4,0) -- (0.8,1.0) -- cycle;
        \node[below, font=\tiny\bfseries, color=AlmVerde] at (0.7,-0.05) {Acutángulo};
        \node[below, font=\tiny, color=AlmGris] at (0.7,-0.25) {todos $<$ 90°};
        % Rectángulo
        \draw[zona]
          (1.9,0) -- (3.3,0) -- (1.9,1.0) -- cycle;
        \draw[rectoang] (1.9,0) -- (2.2,0) -- (2.2,0.3) -- (1.9,0.3);
        \node[below, font=\tiny\bfseries, color=AlmAzulM] at (2.6,-0.05) {Rectángulo};
        \node[below, font=\tiny, color=AlmGris] at (2.6,-0.25) {un ángulo $=90°$};
        % Obtusángulo
        \draw[rectot, fill=AlmTerra!10]
          (4.0,0.5) -- (5.6,0) -- (4.3,1.1) -- cycle;
        \node[below, font=\tiny\bfseries, color=AlmTerra] at (4.9,-0.05) {Obtusángulo};
        \node[below, font=\tiny, color=AlmGris] at (4.9,-0.25) {un ángulo $>90°$};
      \end{tikzpicture}
      \end{center}
    \end{column}
  \end{columns}
  \vspace{0.2cm}
  \begin{formulabox}
    Suma de ángulos de un triángulo $= \alpha + \beta + \gamma = 180°$
    \quad\quad
    Suma de ángulos de un cuadrilátero $= 360°$
  \end{formulabox}
\end{frame}

% ---------------------------------------------------------------
\seccionframe[BAnalizar]{Polígonos en Almería}
% ---------------------------------------------------------------

% ---------------------------------------------------------------
%  FRAME 8 — Almería y polígonos
% ---------------------------------------------------------------
\begin{frame}{Polígonos en la arquitectura almeriense}
  \begin{columns}[t]
    \begin{column}{0.50\textwidth}
      \begin{ejemplobox}
        \small
        \begin{itemize}
          \item \textbf{Catedral de Almería:} fachada con rectángulos,
                arcos y triángulos en las cornisas.
          \item \textbf{Plazas del centro:} cuadriláteros (cuadrados
                y trapecios) en el trazado urbano.
          \item \textbf{Torres de la Alcazaba:} polígonos irregulares
                en la planta de cada torre.
          \item \textbf{Puerto de Almería:} edificios con ventanas
                hexagonales regulares.
        \end{itemize}
      \end{ejemplobox}
    \end{column}
    \begin{column}{0.46\textwidth}
      \begin{center}
      \begin{tikzpicture}[scale=0.72]
        % Silueta estilizada de edificio con formas
        % Base rectangular
        \draw[zona] (0,0) rectangle (3.6,1.8);
        % Ventana rectangular
        \draw[rectov, fill=AlmAzulC!25] (0.3,0.4) rectangle (1.0,1.2);
        \draw[rectov, fill=AlmAzulC!25] (1.4,0.4) rectangle (2.1,1.2);
        \draw[rectov, fill=AlmAzulC!25] (2.5,0.4) rectangle (3.2,1.2);
        % Tejado triangular
        \draw[rectot, fill=AlmTerra!20] (0,1.8) -- (1.8,2.8) -- (3.6,1.8) -- cycle;
        % Hexágono (ventana especial)
        \node[font=\tiny\bfseries, color=AlmAzul] at (1.8,0.2) {Rectángulos};
        \node[font=\tiny\bfseries, color=AlmTerra] at (1.8,2.55) {Triángulo};
      \end{tikzpicture}
      \end{center}
      \begin{notabox}[Dato]
        \scriptsize La Catedral de Almería (s. XVI) es una de las
        pocas catedrales-fortaleza de España: sus muros combinan
        rectángulos y arcos de medio punto.
      \end{notabox}
    \end{column}
  \end{columns}
\end{frame}

% ---------------------------------------------------------------
\seccionframe[BRecordar]{Ejercicios — RECORDAR}
% ---------------------------------------------------------------

% ---------------------------------------------------------------
%  FRAME 9 — Ejercicio RECORDAR
% ---------------------------------------------------------------
\begin{frame}{Ejercicio 1 — RECORDAR \dificultad{1}}
  \begin{recordarbox}
    \begin{enumerate}
      \item ¿Cómo se llama el polígono de 5 lados?
      \item ¿Cuántos vértices tiene un hexágono?
      \item Completa la tabla:

      \vspace{0.2cm}
      \begin{center}
      \small
      \begin{tabular}{|>{\bfseries\color{AlmAzulM}}l|c|c|c|}
        \hline
        \rowcolor{AlmAzul!15}
        \textbf{Figura} & \textbf{Lados} & \textbf{Vértices} &
        \textbf{Suma ángulos} \\
        \hline
        Triángulo  & 3 & \phantom{000} & \phantom{000°} \\
        \hline
        Cuadrilátero & \phantom{0} & 4 & \phantom{000°} \\
        \hline
        Pentágono  & \phantom{0} & \phantom{0} & 540° \\
        \hline
        Hexágono   & \phantom{0} & \phantom{0} & \phantom{000°} \\
        \hline
        Octógono   & \phantom{0} & \phantom{0} & \phantom{000°} \\
        \hline
      \end{tabular}
      \end{center}
    \end{enumerate}
  \end{recordarbox}
  \vspace{0.15cm}
  \begin{notabox}[Pista]
    \small Suma de ángulos $= (n-2)\times 180°$, donde $n$ es el número de lados.
  \end{notabox}
\end{frame}

% ---------------------------------------------------------------
\seccionframe[BComprender]{Ejercicios — COMPRENDER}
% ---------------------------------------------------------------

% ---------------------------------------------------------------
%  FRAME 10 — Ejercicio COMPRENDER
% ---------------------------------------------------------------
\begin{frame}{Ejercicio 2 — COMPRENDER \dificultad{2}}
  \begin{comprenderbox}
    Una plaza de Almería tiene cuatro lados de longitudes:
    \textbf{25 m, 30 m, 25 m y 30 m}. Los cuatro ángulos son de 90°.
    \begin{enumerate}
      \item ¿Qué tipo de cuadrilátero es?
      \item ¿Es regular o irregular? Razona tu respuesta.
      \item ¿Qué condición adicional necesitaría para ser un cuadrado?
    \end{enumerate}
  \end{comprenderbox}
  \vspace{0.2cm}
  \begin{columns}[t]
    \begin{column}{0.48\textwidth}
      \begin{pasobox}[1]
        Identificar el tipo por sus propiedades:
        lados opuestos iguales + 4 ángulos rectos.
      \end{pasobox}
    \end{column}
    \begin{column}{0.48\textwidth}
      \begin{pasobox}[2]
        Comparar con la definición de polígono regular:
        ¿todos los lados son iguales?
      \end{pasobox}
    \end{column}
  \end{columns}
\end{frame}

% ---------------------------------------------------------------
\seccionframe[BAplicar]{Ejercicios — APLICAR}
% ---------------------------------------------------------------

% ---------------------------------------------------------------
%  FRAME 11 — Ejercicio APLICAR
% ---------------------------------------------------------------
\begin{frame}{Ejercicio 3 — APLICAR \dificultad{3}}
  \begin{aplicarbox}
    Identifica los polígonos que aparecen en estos 5 elementos
    de Almería y clasifícalos (regular/irregular):

    \vspace{0.15cm}
    \small
    \begin{tabular}{@{}p{3.2cm}p{3.5cm}p{2.8cm}@{}}
      \toprule
      \textbf{Elemento} & \textbf{Polígono(s)} & \textbf{Reg. / Irreg.} \\
      \midrule
      Fachada Catedral & \phantom{xxxxxxxxxx} & \phantom{xxxxx} \\
      Ventana hexagonal & \phantom{xxxxxxxxxx} & \phantom{xxxxx} \\
      Torre Alcazaba & \phantom{xxxxxxxxxx} & \phantom{xxxxx} \\
      Plaza del Ayuntamiento & \phantom{xxxxxxxxxx} & \phantom{xxxxx} \\
      Panel solar (tejado) & \phantom{xxxxxxxxxx} & \phantom{xxxxx} \\
      \bottomrule
    \end{tabular}
  \end{aplicarbox}
  \vspace{0.15cm}
  \begin{notabox}[Recuerda]
    \small Cuenta los lados visibles. Comprueba si todos los lados
    y ángulos son iguales para determinar si es regular.
  \end{notabox}
\end{frame}

% ---------------------------------------------------------------
\seccionframe[BAnalizar]{Ejercicios — ANALIZAR}
% ---------------------------------------------------------------

% ---------------------------------------------------------------
%  FRAME 12 — Ejercicio ANALIZAR
% ---------------------------------------------------------------
\begin{frame}{Ejercicio 4 — ANALIZAR \dificultad{4}}
  \begin{analizarbox}
    Compara un \textbf{cuadrado} y un \textbf{rectángulo}:
    \begin{enumerate}
      \item ¿Qué tienen en común (número de lados, vértices,
            suma de ángulos)?
      \item ¿En qué se diferencian?
      \item Si ves una figura con 4 ángulos rectos pero sin
            ninguna medida, ¿puedes saber si es cuadrado
            o rectángulo? ¿Por qué?
    \end{enumerate}
  \end{analizarbox}
  \vspace{0.2cm}
  \begin{columns}[c]
    \begin{column}{0.30\textwidth}
      \begin{center}
      \begin{tikzpicture}[scale=0.65]
        \draw[zona] (0,0) rectangle (1.4,1.4);
        \draw[rectoang] (0,0) -- (0.25,0) -- (0.25,0.25) -- (0,0.25);
        \node[below, font=\tiny\bfseries] at (0.7,-0.1) {Cuadrado};
        \node[below, font=\tiny, color=AlmGris] at (0.7,-0.35) {$l = l$};
      \end{tikzpicture}
      \end{center}
    \end{column}
    \begin{column}{0.30\textwidth}
      \begin{center}
      \begin{tikzpicture}[scale=0.65]
        \draw[zona] (0,0) rectangle (2.0,1.1);
        \draw[rectoang] (0,0) -- (0.25,0) -- (0.25,0.25) -- (0,0.25);
        \node[below, font=\tiny\bfseries] at (1.0,-0.1) {Rectángulo};
        \node[below, font=\tiny, color=AlmGris] at (1.0,-0.35) {$l \neq w$};
      \end{tikzpicture}
      \end{center}
    \end{column}
    \begin{column}{0.34\textwidth}
      \small\color{AlmAzulM}
      Ambos: 4 lados, 4 vértices, 360°. \\[4pt]
      Diferencia: ¿todos los lados iguales?
    \end{column}
  \end{columns}
\end{frame}

% ---------------------------------------------------------------
%  FRAME 13 — Error frecuente
% ---------------------------------------------------------------
\begin{frame}{¡Cuidado con estos errores!}
  \begin{errorbox}
    \begin{itemize}
      \item \incorrecto\ \textbf{«El círculo es un polígono»}
            \\
            \correcto\ Un círculo tiene un contorno \emph{curvo}, no lados rectos.
            \textbf{No es un polígono.}
      \vspace{0.3cm}
      \item \incorrecto\ \textbf{«El hexágono tiene 5 lados»}
            \\
            \correcto\ HEXA = 6. \quad
            \textbf{HEXÁgono $\rightarrow$ 6 lados}.
            \\
            Pentágono (PENTA $= 5$) es el de 5 lados.
      \vspace{0.3cm}
      \item \incorrecto\ \textbf{«Todo cuadrilátero es un cuadrado»}
            \\
            \correcto\ Solo si los 4 lados son iguales
            \textbf{y} los 4 ángulos son rectos.
    \end{itemize}
  \end{errorbox}
\end{frame}

% ---------------------------------------------------------------
%  FRAME 14 — Resumen
% ---------------------------------------------------------------
\begin{frame}{Resumen de la sesión}
  \begin{resumenbox}
    \begin{itemize}
      \item Un \textbf{polígono} es una figura plana cerrada con lados rectos.
      \item Se clasifica por número de lados: triángulo (3) hasta octógono (8).
      \item En todo polígono: \#lados $=$ \#vértices $=$ \#ángulos.
      \item \textbf{Regular}: todos los lados y ángulos iguales.
            \textbf{Irregular}: no se cumple.
      \item Triángulos por lados: equilátero, isósceles, escaleno.
      \item Triángulos por ángulos: acutángulo, rectángulo, obtusángulo.
      \item Suma de ángulos: triángulo $= 180°$, cuadrilátero $= 360°$,
            en general $(n-2)\times 180°$.
    \end{itemize}
  \end{resumenbox}
\end{frame}

% ---------------------------------------------------------------
%  FRAME 15 — Evidencia del día
% ---------------------------------------------------------------
\begin{frame}{Evidencia del día}
  \begin{evidenciabox}
    \textbf{Catálogo de figuras planas de Almería}

    \vspace{0.3cm}
    Elabora un catálogo con al menos \textbf{5 edificios o elementos}
    de Almería. Para cada uno indica:

    \vspace{0.15cm}
    \begin{enumerate}
      \item Nombre del edificio o elemento.
      \item Polígono(s) identificado(s).
      \item Número de lados y vértices.
      \item Regular o irregular — justifica con las propiedades.
      \item Dibuja el polígono (a mano o con regla y compás).
    \end{enumerate}

    \vspace{0.2cm}
    \textbf{Criterio de éxito:} los 5 polígonos están correctamente
    identificados, clasificados y dibujados.
  \end{evidenciabox}
\end{frame}

\end{document}
